\documentclass{article}%
\usepackage[T1]{fontenc}%
\usepackage[utf8]{inputenc}%
\usepackage{lmodern}%
\usepackage{textcomp}%
\usepackage{lastpage}%
\usepackage{geometry}%
\geometry{margin=1in,headheight=14pt,headsep=25pt}%
\usepackage{hyperref}%
\usepackage{enumitem}%
\usepackage{fancyhdr}%
\usepackage{xcolor}%
\usepackage{url}%
\usepackage{breakurl}%
\usepackage{amsmath}%
\usepackage{amssymb}%
\usepackage{mathtools}%
\usepackage{amsthm}%
\usepackage{thmtools}%
\usepackage{float}%
\usepackage{subcaption}%
\usepackage[font=small,labelfont=bf]{caption}%
%

            \hypersetup{
                colorlinks=true,
                linkcolor=blue,
                filecolor=magenta,
                urlcolor=blue,
                breaklinks=true,
                bookmarks=true,
                pdfborder={0 0 0}
            }
            \urlstyle{same}
            \Urlmuskip=0mu plus 1mu  % URL breaking in nicer spots

            % Math configuration
            \everymath{\displaystyle}
            \setlength{\jot}{10pt}  % Increase spacing between equations

            \pagestyle{fancy}
            \fancyhf{}
            \rhead{Generated on \today}
            \cfoot{\thepage}

            % Better Unicode support
            \usepackage[utf8]{inputenc}
            \usepackage{newunicodechar}
            \setlength{\itemsep}{1em} % Adds spacing between items

            % Configure itemize settings globally
            \setlist[itemize]{nosep}
            \setlist[itemize]{leftmargin=*}

            % Remove section numbering
            \renewcommand{\thesection}{}
            \renewcommand{\thesubsection}{}

            % Adjust caption spacing around figures
            \setlength{\abovecaptionskip}{0pt}
            \setlength{\belowcaptionskip}{5pt}
        %
\lhead{Lecture: 16L2{-}ethics}%
\title{16L2{-}ethics}%
\author{Generated by Scribe.AI}%
\date{\today}%
%
\begin{document}%
\normalsize%
\maketitle%
\section{Slide 1}%
\label{sec:Slide1}%
\subsection{Content}%
%
\subsection{Description}%
This slide is a title slide for a lecture on Introduction to Data Science, likely part of a course called Purdue CS24200.  It displays the course title, the course number, and the instructor's name (Tony Bergstrom).  A footer indicates that the slides are adapted from another source, specifically Jen Neville's CS 24200 course.  The overall design is simple, using a dark gray background and light gray text, making it visually clear and professional.%
\subsection{Figures}%
\begin{figure}[H]%
\centering%
\end{figure}

%
\newpage%
\section{Slide 2}%
\label{sec:Slide2}%
\subsection{Content}%

\textbf{Bias and Fairness in Data}
%
\subsection{Description}%
This slide, likely part of a lecture on Artificial Intelligence Basics, displays a single heading: "Bias and Fairness in Data."  It's a simple title slide, indicating the topic of the subsequent content, which will likely discuss issues of bias and fairness in datasets used for AI applications.%
\subsection{Figures}%
\begin{figure}[H]%
\centering%
\end{figure}

%
\newpage%
\section{Slide 3}%
\label{sec:Slide3}%
\subsection{Content}%

\textbf{Why data analysis methods are helpful...} \\
Provide algorithms for tasks not previously amenable to automation (e.g. image analysis) \\
Provide advantages over humans doing similar tasks: \\
\begin{itemize}
    \item Inexpensive/scalable/fast
    \item Consistent and verifiable
    \item More accurate*
    \item More fair*, less subject to social biases
\end{itemize}
\textit{*Maybe, sometimes} \\
Source: Lawrence Hunter, Colorado School of Mines
%
\subsection{Description}%
This slide, likely from an Artificial Intelligence Basics course, discusses the benefits of data analysis methods.  It lists several advantages of using algorithms over human analysis for tasks like image analysis.  The bullet points highlight the key benefits:  inexpensive/scalable/fast processing, consistent and verifiable results, increased accuracy (*with a caveat*), and reduced susceptibility to social biases (*with a caveat*).  The asterisk next to "More accurate" and "More fair" suggests that these advantages are not absolute and may depend on the specific application or dataset. The slide's source is attributed to Lawrence Hunter from the Colorado School of Mines.%
\subsection{Figures}%
\begin{figure}[H]%
\centering%
\end{figure}

%
\newpage%
\section{Slide 4}%
\label{sec:Slide4}%
\subsection{Content}%
%
\subsection{Description}%
This slide, likely from an Artificial Intelligence Basics course, presents a visual analysis of funding and mergers & acquisitions (M&A) activity related to AI startups.  It displays two key graphs. The first graph, on the left, shows the amount of funding (in millions of dollars) to AI startups over time, from 2013 to 2019, broken down by quarter.  The data points are represented by bars, and the overall trend shows a significant increase in funding over the period.  The title indicates that Q2 2019 saw a record amount of funding, reaching \$7.4 billion.  The graph includes a subtitle indicating that the data is from Q1 2013 to Q2 2019, and that additional data can be accessed by swiping right. The source of the data is credited to CB Insights. The second graph, on the right, shows the number of M&A deals in the AI sector from 2010 to 2019.  The data points are also represented by bars, and the graph shows fluctuations in the number of deals over time.  The graph includes a note indicating that the data for 2017 is up to December 4th. The overall design of the slide is to visually represent the growth of the AI industry in terms of funding and M&A activity. The context of the course is crucial for understanding the significance of these trends in the development and adoption of AI technologies.%
\subsection{Figures}%
\begin{figure}[H]%
\centering%
\begin{subfigure}[b]{0.15\textwidth}%
\centering%
\includegraphics[width=\textwidth]{/Users/ashoksaravanan/Coding/ScribeLec/server/output/CS 243/lectures/16L2-ethics/figures/4.1.png}%
\caption{}%
\end{subfigure}%
\hspace{0.05\textwidth}%
\begin{subfigure}[b]{0.15\textwidth}%
\centering%
\includegraphics[width=\textwidth]{/Users/ashoksaravanan/Coding/ScribeLec/server/output/CS 243/lectures/16L2-ethics/figures/4.2.png}%
\caption{}%
\end{subfigure}%
\hspace{0.05\textwidth}%
\begin{subfigure}[b]{0.15\textwidth}%
\centering%
\includegraphics[width=\textwidth]{/Users/ashoksaravanan/Coding/ScribeLec/server/output/CS 243/lectures/16L2-ethics/figures/4.3.png}%
\caption{}%
\end{subfigure}%
\end{figure}

%
\newpage%
\section{Slide 5}%
\label{sec:Slide5}%
\subsection{Content}%

\textbf{Examples in Society} \\
\begin{itemize}
    \item Web advertising
    \item Lending
    \item Ridesharing
    \item Public policy
\end{itemize}
\textit{Source: Richard Johansson, University of Gothenburg}
%
\subsection{Description}%
This slide, likely from an Artificial Intelligence Basics course, presents examples of how AI is used in various societal contexts.  It displays four distinct examples: web advertising, lending, ridesharing, and public policy.  Each example is visually represented by a small image or graphic.  For example, the "Web advertising" section likely shows an image of a collage of various images, suggesting the use of AI in targeted advertising.  The "Lending" section likely shows a graphic depicting a computer or robot with a balance scale, illustrating AI's role in loan applications and risk assessment.  The "Ridesharing" section likely shows a map interface, highlighting AI's use in ride-hailing services.  The "Public policy" section likely shows a graphic related to policy-making, suggesting AI's potential in analyzing and predicting trends.  The slide's title, "Examples in Society," indicates that the focus is on the practical applications of AI in real-world scenarios.  The source is credited to Richard Johansson from the University of Gothenburg.%
\subsection{Figures}%
\begin{figure}[H]%
\centering%
\begin{subfigure}[b]{0.15\textwidth}%
\centering%
\includegraphics[width=\textwidth]{/Users/ashoksaravanan/Coding/ScribeLec/server/output/CS 243/lectures/16L2-ethics/figures/5.1.png}%
\caption{}%
\end{subfigure}%
\end{figure}

%
\newpage%
\section{Slide 6}%
\label{sec:Slide6}%
\subsection{Content}%

\textbf{ChatGPT Is the Most Tried} \\
\textbf{AI Tool And Users Stick to It} \\
Share of U.S. respondents who have used \\
the following AI tools in the past 12 months \\
\begin{itemize}
    \item ChatGPT \quad 20.9\% \quad 89\%
    \item Jasper Chat \quad 12.9\% \quad 69\%
    \item GetGenie \quad 11.3\% \quad 66\%
    \item Simplified \quad 10.9\% \quad 69\%
    \item DALL-E \quad 10.1\% \quad 72\%
    \item Midjourney \quad 10.0\% \quad 65\%
    \item YouChat \quad 9.8\% \quad 80\%
    \item DeepL \quad 9.6\% \quad 72\%
\end{itemize}
\textit{1,237 U.S. respondents (18-64 y/o) surveyed Mar.-Apr. 2023} \\
\textit{Source: Statista Consumer Insights} \\
\textbf{ChatGPT Sprints to} \\
\textbf{One Million Users} \\
Time it took for selected online services \\
to reach one million users \\
\begin{itemize}
    \item Netflix \quad 1999 \quad 3.5 years
    \item Kickstarter \quad 2009 \quad 2.5 years
    \item Airbnb \quad 2008 \quad 2.5 years
    \item Twitter \quad 2006 \quad 2 years
    \item Foursquare \quad 2009 \quad 13 months
    \item Facebook \quad 2004 \quad 10 months
    \item Dropbox \quad 2008 \quad 7 months
    \item Spotify \quad 2008 \quad 5 months
    \item Instagram \quad 2010 \quad 2.5 months
\end{itemize}
\textit{ChatGPT 2022 \quad 5 days} \\
\textit{* one million backers ** one million nights booked *** one million downloads} \\
\textit{Source: Company announcements via Business Insider/LinkedIn}
%
\subsection{Description}%
This slide from an Artificial Intelligence Basics course presents two charts related to ChatGPT and other AI tools. The first chart, "ChatGPT Is the Most Tried AI Tool And Users Stick to It," displays the percentage of U.S. respondents who have used various AI tools in the past 12 months.  Each tool (ChatGPT, Jasper Chat, GetGenie, Simplified, DALL-E, Midjourney, YouChat, and DeepL) is listed with its corresponding usage percentage.  A second percentage, likely representing the percentage of users who would use the tool again, is also provided for each tool.  The second chart, "ChatGPT Sprints to One Million Users," shows the time it took for various online services to reach one million users.  Each service (Netflix, Kickstarter, Airbnb, Twitter, Foursquare, Facebook, Dropbox, Spotify, and Instagram) is listed with the year it was launched and the time it took to reach one million users.  The data is presented in a bar graph format, making it easy to compare the usage and growth rates of different AI tools and online services.  The source of the data is credited to Statista.  The context of the course is about understanding the adoption and growth rates of AI tools and services.%
\subsection{Figures}%
\begin{figure}[H]%
\centering%
\begin{subfigure}[b]{0.15\textwidth}%
\centering%
\includegraphics[width=\textwidth]{/Users/ashoksaravanan/Coding/ScribeLec/server/output/CS 243/lectures/16L2-ethics/figures/6.1.png}%
\caption{}%
\end{subfigure}%
\hspace{0.05\textwidth}%
\begin{subfigure}[b]{0.15\textwidth}%
\centering%
\includegraphics[width=\textwidth]{/Users/ashoksaravanan/Coding/ScribeLec/server/output/CS 243/lectures/16L2-ethics/figures/6.2.png}%
\caption{}%
\end{subfigure}%
\end{figure}

%
\newpage%
\section{Slide 7}%
\label{sec:Slide7}%
\subsection{Content}%

\textbf{What could be the problem?}
%
\subsection{Description}%
This slide, likely part of a presentation or lecture on Artificial Intelligence Basics, poses a question: "What could be the problem?".  The slide contains only this question, written in a bold, sans-serif font.  There are no images, graphs, or other content.  The context suggests that the slide is designed to prompt discussion or introduce a problem-solving exercise related to AI.  The question is likely intended to encourage students to consider potential challenges or limitations in AI systems or applications.%
\subsection{Figures}%
\begin{figure}[H]%
\centering%
\end{figure}

%
\newpage%
\section{Slide 8}%
\label{sec:Slide8}%
\subsection{Content}%

\textbf{What could be the problem?} \\
Data are (about) people, can cause harm \\
Algorithmic outcomes often not explainable \\
A lot of data incidentally produced by daily life: \\
\begin{itemize}
    \item Social media
    \item Ubiquitous cameras, microphones, location tracking
    \item Medical treatment
    \item Public records
\end{itemize} \\
Important new uses \\
\begin{itemize}
    \item Legal: Surveillance, "predictive" policing, sentencing, fraud detection, military applications
    \item Economic: School admissions, hiring/promotion, loans, insurance, accounting controls, advertising
    \item Medical: Health insurance, diagnosis, decision making
    \item Artistic: Copyright/Ownership, Generating new work/performance
\end{itemize}
%
\subsection{Description}%
This slide, likely from an Artificial Intelligence Basics course, presents a list of potential problems and uses of AI.  The first part, "What could be the problem?", highlights concerns about AI systems.  It lists that data about people can cause harm, algorithmic outcomes are often not explainable, and a large amount of data is incidentally collected in daily life through various sources like social media, ubiquitous surveillance devices, medical records, and public records.  The second part, "Important new uses," lists various applications of AI in different sectors.  These include legal applications like surveillance and predictive policing, economic applications like school admissions and loan decisions, medical applications like health insurance and diagnosis, and artistic applications like copyright and generating new works.  The slide is designed to initiate a discussion about the ethical and societal implications of AI, prompting consideration of both the benefits and potential harms of AI technologies.%
\subsection{Figures}%
\begin{figure}[H]%
\centering%
\end{figure}

%
\newpage%
\section{Slide 9}%
\label{sec:Slide9}%
\subsection{Content}%

\textbf{Ethical concerns} \\
\textbf{Preserving privacy} \\
\begin{itemize}
    \item Methods for handling sensitive data
    \item Uses of data science that undermine privacy
\end{itemize} \\
\textbf{Avoiding bias} \\
\begin{itemize}
    \item Data selection and unintentional red-lining
    \item Re-inscription of existing biases
\end{itemize} \\
\textbf{Mitigating malicious attacks} \\
\begin{itemize}
    \item Intentional subversion of machine learning systems
    \item Hazards of learning from the open internet
\end{itemize} \\
Source: Lawrence Hunter, Colorado School of Mines
%
\subsection{Description}%
This slide, likely from an Artificial Intelligence Basics course, lists ethical concerns related to AI.  It's organized into three main categories: Preserving privacy, Avoiding bias, and Mitigating malicious attacks.  Under each category, specific issues are outlined.  For example, under "Preserving privacy," the slide mentions methods for handling sensitive data and how data science itself can sometimes undermine privacy.  Under "Avoiding bias," the slide highlights data selection issues, unintentional biases (like red-lining), and the re-inscription of existing biases in datasets.  Finally, under "Mitigating malicious attacks," the slide points out the potential for intentional subversion of AI systems and the risks of learning from publicly available data on the internet.  The slide is designed to raise awareness of the ethical considerations that must be addressed when developing and deploying AI systems.  The source is credited to Lawrence Hunter, Colorado School of Mines.%
\subsection{Figures}%
\begin{figure}[H]%
\centering%
\end{figure}

%
\newpage%
\section{Slide 10}%
\label{sec:Slide10}%
\subsection{Content}%

\textbf{Anonymized data isn't always safe} \\
In 1997, Latanya Sweeney identified the Governor of Massachusetts' (Weld) medical records. \\
\begin{itemize}
    \item Massachusetts released hospital records anonymized by removing names, addresses and SSNs
    \item Voter records have name, address, ZIP code, birth date, and sex of every voter
    \item Sweeney used zip code, birthdate and gender to uniquely identify Weld's records
    \item 87\% of US identified by zip, birthdate \& gender
\end{itemize} \\
Similar with Netflix (using IMDB) and search logs \\
Source: Lawrence Hunter, Colorado School of Mines
%
\subsection{Description}%
This slide, likely from an Artificial Intelligence Basics course, discusses the potential pitfalls of anonymized data.  It highlights a case study from 1997 where Latanya Sweeney, using publicly available anonymized data, was able to identify the Governor of Massachusetts' medical records.  The slide lists the steps involved: Massachusetts released anonymized hospital records, voter records contained detailed personal information, and Sweeney used readily available data (zip code, birthdate, and gender) to re-identify the Governor.  The slide also emphasizes that a similar issue exists with other anonymized datasets, such as Netflix's IMDB data and search logs.  The source is credited to Lawrence Hunter, Colorado School of Mines.  The overall message is that seemingly anonymized data can still be used to identify individuals, highlighting the importance of careful data handling and privacy considerations in AI applications.%
\subsection{Figures}%
\begin{figure}[H]%
\centering%
\end{figure}

%
\newpage%
\section{Slide 11}%
\label{sec:Slide11}%
\subsection{Content}%

\textbf{Privacy != Security} \\
Some data cannot be anonymized \\
\begin{itemize}
    \item Genome sequences are inherently identifying
    \item Even a few hundred well-picked SNPs...
\end{itemize} \\
Often, people's desires about their data involve questions of trust \\
\begin{itemize}
    \item Willingness to share medical data with academic researchers, but not pharmaceutical companies
\end{itemize} \\
Privacy is not a binary value \\
\begin{itemize}
    \item Different sorts of exposure to different sorts of people evoke different responses
\end{itemize} \\
Source: Lawrence Hunter, Colorado School of Mines
%
\subsection{Description}%
This slide, likely from an Artificial Intelligence Basics course, discusses the complexities of privacy and security in the context of data handling.  It argues that privacy and security are not equivalent and that some data inherently cannot be anonymized.  The slide specifically points out that genome sequences and even a small number of carefully selected genetic markers (SNPs) are inherently identifying.  It further emphasizes that people's willingness to share data is often nuanced and depends on the context and the entity requesting the data.  For example, people might be willing to share medical data with academic researchers but not with pharmaceutical companies.  The slide concludes by stating that privacy is not a simple binary concept; different levels of exposure to different types of people and situations can lead to different privacy concerns and responses.  The source is attributed to Lawrence Hunter, from the Colorado School of Mines.%
\subsection{Figures}%
\begin{figure}[H]%
\centering%
\end{figure}

%
\newpage%
\section{Slide 12}%
\label{sec:Slide12}%
\subsection{Content}%

\textbf{Bias in Data Models} \\
"Objective" algorithms are thought to be free of the biases that plague people. \\
Algorithms, especially ones that learn, can inadvertently re-inscribe those biases \\
Algorithms are opaque, hard to interrogate \\
Increasingly widespread \\
Source: Lawrence Hunter, Colorado School of Mines
%
\subsection{Description}%
This slide, likely from an Artificial Intelligence Basics course, discusses the issue of bias in data models.  It states that "objective" algorithms are often assumed to be free from the biases present in human-created data. However, the slide points out that algorithms, especially those that learn from data, can unintentionally perpetuate or even amplify existing biases.  The slide further notes that these algorithms are often opaque, making it difficult to understand how they arrive at their conclusions and identify potential biases.  Finally, the slide emphasizes the increasing prevalence of these algorithms in various applications. The source is credited to Lawrence Hunter, from the Colorado School of Mines.%
\subsection{Figures}%
\begin{figure}[H]%
\centering%
\end{figure}

%
\newpage%
\section{Slide 13}%
\label{sec:Slide13}%
\subsection{Content}%

\textbf{Proxies for discrimination} \\
Illegal, and generally perceived as wrong to make choices based on race, gender, religion, national origin, etc. \\
However, proxies for these are everywhere: \\
\begin{itemize}
    \item Zip codes (gender, race, national origin)
    \item Names (gender, race, national origin)
    \item Purchase histories (including movies or tv shows)
\end{itemize} \\
Machine learning that uses biased historical record + any proxy is likely to re-inscribe bias \\
Source: Lawrence Hunter, Colorado School of Mines
%
\subsection{Description}%
This slide, likely from an Artificial Intelligence Basics course, discusses the issue of proxies for discriminatory practices in data.  It begins by stating that making choices based on race, gender, religion, or national origin is illegal and generally considered wrong.  However, the slide highlights that proxies for these factors are frequently present in data.  The examples given are: zip codes, which can correlate with demographic information; names, which can reveal gender, race, or national origin; and purchase histories, which can reflect socioeconomic status or other sensitive attributes.  The slide concludes by stating that machine learning models trained on data containing these proxies, combined with existing biased historical records, are likely to perpetuate and even amplify those biases.  The source is attributed to Lawrence Hunter, from the Colorado School of Mines.%
\subsection{Figures}%
\begin{figure}[H]%
\centering%
\end{figure}

%
\newpage%
\section{Slide 14}%
\label{sec:Slide14}%
\subsection{Content}%

\textbf{Discrimination in online ad delivery} \\
Sweeney observed in 2013 that black-identifying names turned out to be much more \\
likely than white-identifying names to generate ads that including the word "arrest" (60 \\
per cent versus 48 per cent). \\
Google uses a learning algorithm to place ads that are most often clicked on. \\
Likely to be a reflection of people clicking on those ads more for 'black' names \\
Source: Lawrence Hunter, Colorado School of Mines
%
\subsection{Description}%
This slide from an Artificial Intelligence Basics course discusses a case study of potential bias in online advertising.  It describes an observation by Sweeney in 2013 that black-identifying names were more likely to receive ads containing the word "arrest" than white-identifying names (60% vs. 48%).  The slide suggests that Google's algorithm, which prioritizes ads based on click-through rates, might be amplifying this bias.  It posits that the higher click-through rate for ads shown to black-identifying names could be a reflection of users clicking on those ads more frequently, rather than a systematic bias in the algorithm itself. The source is attributed to Lawrence Hunter from the Colorado School of Mines.%
\subsection{Figures}%
\begin{figure}[H]%
\centering%
\end{figure}

%
\newpage%
\section{Slide 15}%
\label{sec:Slide15}%
\subsection{Content}%

\textbf{Women less likely to be shown ads for high-paid jobs on Google, study shows} \\
Automated testing and analysis of company's advertising system reveals male job seekers are shown far more adverts for high-paying executive jobs \\
\textbf{Female job seekers are much less likely to be shown adverts on Google for highly paid jobs than men, researchers have found.} \\
The team of researchers from Carnegie Mellon built an automated testing rig called AdFisher that pretended to be a series of male and female job seekers. Their 17,370 fake profiles only visited jobseeker sites and were shown 600,000 adverts which the team tracked and analysed. \\
\textbf{One experiment showed that Google displayed adverts for a career coaching service for executive jobs 1,852 times to the male group and only 318 times to the female group.}
%
\subsection{Description}%
This slide, likely from a presentation or news article, discusses a study revealing potential gender bias in Google's advertising system.  The study, conducted by researchers from Carnegie Mellon University, used an automated testing rig (called AdFisher) to simulate job seekers.  17,370 fake profiles, representing both male and female job seekers, were created.  These profiles visited job-seeking websites and were exposed to various advertisements.  The results showed that male job seekers were significantly more likely (1,852 times) to see ads for high-paying executive jobs compared to female job seekers (318 times).  This suggests a potential bias in Google's ad delivery system, potentially favoring male job seekers.  The slide highlights the methodology of the study, emphasizing the automated testing rig and the large number of simulated profiles used to gather data.  The context of the course (Artificial Intelligence Basics) would likely focus on the implications of such bias in algorithms and the potential for unintended consequences in automated systems.%
\subsection{Figures}%
\begin{figure}[H]%
\centering%
\begin{subfigure}[b]{0.15\textwidth}%
\centering%
\includegraphics[width=\textwidth]{/Users/ashoksaravanan/Coding/ScribeLec/server/output/CS 243/lectures/16L2-ethics/figures/15.1.png}%
\caption{}%
\end{subfigure}%
\end{figure}

%
\newpage%
\section{Slide 16}%
\label{sec:Slide16}%
\subsection{Content}%

\textbf{Machine Bias} \\
There's software used across the country to predict future criminals. And it's biased against blacks. \\
by Julia Angwin, Jeff Larson, Surya Mattu and Lauren Kirchner, ProPublica \\
May 23, 2016
%
\subsection{Description}%
This slide, likely from an Artificial Intelligence Basics course, presents a headline about a ProPublica article.  The title "Machine Bias" indicates the focus on potential bias in machine learning algorithms.  The text states that predictive software used for criminal justice is biased against black individuals.  The article's authors, Julia Angwin, Jeff Larson, Surya Mattu, and Lauren Kirchner from ProPublica, are credited.  The date of the article, May 23, 2016, is also included.  The slide likely serves as an introduction to a discussion about algorithmic bias in the context of AI, highlighting a real-world example of how AI systems can perpetuate existing societal biases.%
\subsection{Figures}%
\begin{figure}[H]%
\centering%
\begin{subfigure}[b]{0.15\textwidth}%
\centering%
\includegraphics[width=\textwidth]{/Users/ashoksaravanan/Coding/ScribeLec/server/output/CS 243/lectures/16L2-ethics/figures/16.1.png}%
\caption{}%
\end{subfigure}%
\hspace{0.05\textwidth}%
\begin{subfigure}[b]{0.15\textwidth}%
\centering%

\includegraphics[width=\textwidth]{/Users/ashoksaravanan/Coding/ScribeLec/server/output/CS 243/lectures/16L2-ethics/figures/16.2.png}%
\caption{}%
\end{subfigure}%
\end{figure}

%
\newpage%
\section{Slide 17}%
\label{sec:Slide17}%
\subsection{Content}%
%
\subsection{Description}%
This slide, likely from an Artificial Intelligence Basics course, presents a visual comparison of the distribution of "decile scores" for black and white defendants.  Two bar charts are shown side-by-side.  The chart on the left displays the distribution of decile scores for black defendants, and the chart on the right displays the distribution for white defendants.  The x-axis of each chart represents the decile score (from 1 to 10), and the y-axis represents the count of defendants falling into each decile.  The charts visually illustrate the frequency of defendants in each decile score category for both groups.  The source is attributed to Richard Johansson from the University of Gothenburg.  The context of the course suggests that this slide is likely part of a discussion about potential bias in algorithms, specifically in the context of risk assessment systems like the COMPAS (Correctional Offender Management Profiling for Alternative Sanctions) system, which is often used to predict recidivism. The visual comparison of the distributions is intended to highlight potential disparities in how the system assigns scores to different racial groups.%
\subsection{Figures}%
\begin{figure}[H]%
\centering%
\begin{subfigure}[b]{0.15\textwidth}%
\centering%
\includegraphics[width=\textwidth]{/Users/ashoksaravanan/Coding/ScribeLec/server/output/CS 243/lectures/16L2-ethics/figures/17.1.png}%
\caption{}%
\end{subfigure}%
\end{figure}

%
\newpage%
\section{Slide 18}%
\label{sec:Slide18}%
\subsection{Content}%

\textbf{Can Predictive Systems Discriminate Unfairly} \\
When predictive systems are more widely applied in situations where they affect people's \\
lives, they may run into troubles concerning \textbf{anti-discrimination} laws \\
There is no consensus about how to \textbf{define} and even less how to \textbf{deal} with the problem \\
Source: Richard Johansson, University of Gothenburg
%
\subsection{Description}%
This slide, likely from an Artificial Intelligence Basics course, discusses the potential for bias and unfair discrimination in predictive systems.  It highlights the growing concern that as predictive systems become more prevalent in areas affecting people's lives, they may violate anti-discrimination laws.  The slide emphasizes the lack of consensus on how to define and address this problem.  The source is attributed to Richard Johansson from the University of Gothenburg, suggesting this is a recognized area of concern in the field of AI. The context of the course is focused on the ethical and societal implications of AI, particularly the potential for bias and discrimination in algorithms.%
\subsection{Figures}%
\begin{figure}[H]%
\centering%
\begin{subfigure}[b]{0.15\textwidth}%
\centering%
\includegraphics[width=\textwidth]{/Users/ashoksaravanan/Coding/ScribeLec/server/output/CS 243/lectures/16L2-ethics/figures/18.1.png}%
\caption{}%
\end{subfigure}%
\end{figure}

%
\newpage%
\section{Slide 19}%
\label{sec:Slide19}%
\subsection{Content}%

\textbf{The Legal Situation} \\
Anti-discrimination laws in several countries prohibit unfair treatment of people based on protected \\
attributes (e.g. gender, race) \\
These laws often evaluate the fairness of a decision making process by means of two distinct notions: \\
\begin{itemize}
    \item Disparate treatment: if it's decisions are (partly) based on the subject's protected attribute \\
    information \\
    \item Disparate impact: if it's outcomes disproportionally hurt (or benefit) people with certain sensitive \\
    attribute values
\end{itemize}
Are these ideas in conflict? \\
How can they be operationalized? \\
Source: Richard Johansson, University of Gothenburg
%
\subsection{Description}%
This slide, likely from an Artificial Intelligence Basics course, discusses the legal aspects of fairness in algorithms.  It introduces the concept of anti-discrimination laws in various countries that prohibit unfair treatment based on protected attributes like gender and race.  The slide outlines two key legal notions used to evaluate the fairness of a decision-making process: disparate treatment and disparate impact.  Disparate treatment refers to situations where decisions are (partially) based on a protected attribute.  Disparate impact refers to situations where the outcomes of a process disproportionately harm or benefit people with certain sensitive attributes.  The slide concludes by posing questions about the potential conflict between these two concepts and how they can be practically applied (operationalized) in the context of AI systems.  The source is attributed to Richard Johansson from the University of Gothenburg, indicating an academic perspective on the topic.%
\subsection{Figures}%
\begin{figure}[H]%
\centering%
\begin{subfigure}[b]{0.15\textwidth}%
\centering%
\includegraphics[width=\textwidth]{/Users/ashoksaravanan/Coding/ScribeLec/server/output/CS 243/lectures/16L2-ethics/figures/19.1.png}%
\caption{}%
\end{subfigure}%
\end{figure}

%
\newpage%
\section{Slide 20}%
\label{sec:Slide20}%
\subsection{Content}%

\textbf{Example: The Compass System} \\
\textbf{Prediction Fails Differently for Black Defendants} \\
\begin{tabular}{ | l | l | l | } 
 \hline
 & WHITE & AFRICAN AMERICAN \\ 
 \hline
 Labeled Higher Risk, But Didn't Re-Offend & 23.5\% & 44.9\% \\ 
 \hline
 Labeled Lower Risk, Yet Did Re-Offend & 47.7\% & 28.0\% \\ 
 \hline
\end{tabular} \\
Non-offending blacks get higher scores \\
• NB: race is not a feature in the model \\
So the recall for the “positive” case isn’t equal for the groups \\
Source: Richard Johansson, University of Gothenburg
%
\subsection{Description}%
This slide, from an Artificial Intelligence Basics course, presents an example of potential bias in the COMPAS (Correctional Offender Management Profiling for Alternative Sanctions) system.  It shows a table comparing the results of the COMPAS risk assessment for white and African American defendants.  The table displays the percentages of defendants who were labeled as higher risk but did not re-offend, and those who were labeled as lower risk but did re-offend.  The percentages are significantly different between the two groups.  Critically, the slide highlights that "race is not a feature in the model," meaning the algorithm is not explicitly using race as an input.  This demonstrates how a system, even without explicitly using race as a factor, can still produce disparate outcomes based on race.  The conclusion is that the recall for the "positive" case (i.e., identifying those at high risk) is not equal for the two groups.  This is a key example of how seemingly unbiased algorithms can exhibit bias in their outputs, potentially leading to unfair or discriminatory outcomes. The source is attributed to Richard Johansson from the University of Gothenburg, suggesting an academic context for the example.%
\subsection{Figures}%
\begin{figure}[H]%
\centering%
\begin{subfigure}[b]{0.15\textwidth}%
\centering%
\includegraphics[width=\textwidth]{/Users/ashoksaravanan/Coding/ScribeLec/server/output/CS 243/lectures/16L2-ethics/figures/20.1.png}%
\caption{}%
\end{subfigure}%
\end{figure}

%
\newpage%
\section{Slide 21}%
\label{sec:Slide21}%
\subsection{Content}%

\textbf{When Machine Learning Reinforces Stereotypes} \\
ML model can pick up subtle biases from the training set \\
Including various stereotypes \\
\begin{itemize}
    \item Images, text, ...
\end{itemize} \\
Source: Richard Johansson, University of Gothenburg
%
\subsection{Description}%
This slide, likely from an Artificial Intelligence Basics course, discusses the potential for bias in machine learning models.  It states that machine learning (ML) models can inadvertently learn and perpetuate stereotypes from the training data.  The slide specifically mentions that biases can be present in various forms of data, including images and text.  The source is attributed to Richard Johansson from the University of Gothenburg, suggesting an academic perspective on the topic of algorithmic bias.%
\subsection{Figures}%
\begin{figure}[H]%
\centering%
\begin{subfigure}[b]{0.15\textwidth}%
\centering%
\includegraphics[width=\textwidth]{/Users/ashoksaravanan/Coding/ScribeLec/server/output/CS 243/lectures/16L2-ethics/figures/21.1.png}%
\caption{}%
\end{subfigure}%
\end{figure}

%
\newpage%
\section{Slide 22}%
\label{sec:Slide22}%
\subsection{Content}%
%
\subsection{Description}%
This slide, titled "Collecting Training Data," presents a visual representation of the process of training a machine learning model.  It uses a flowchart-like diagram with two distinct paths.  The top path depicts a scenario where a predictor is built using machine learning on a sample of a population.  The text below the predictor node indicates that this predictor is "Built to make accurate predictions on the Population samples."  The bottom path shows a similar process, but with the crucial difference that the resulting predictor is "But not predictive of the True Population."  This implies that the model trained on the sample data is not generalizable to the entire population.  The diagram highlights the importance of representative sampling and the potential for overfitting when training machine learning models.  The diagram shows two separate paths, each starting with "Population Samples" and ending with a "Predictor."  The top path's predictor is accurate for the sample, while the bottom path's predictor is not accurate for the true population.  The diagram is part of a lecture on Artificial Intelligence Basics, likely focusing on the practical aspects of machine learning and the importance of data quality and representation.%
\subsection{Figures}%
\begin{figure}[H]%
\centering%
\end{figure}

%
\newpage%
\section{Slide 23}%
\label{sec:Slide23}%
\subsection{Content}%

\textbf{Amazon scraps secret AI recruiting tool that showed bias against women}
%
\subsection{Description}%
This slide, likely from an Artificial Intelligence Basics course, presents a news article excerpt about Amazon's AI recruiting tool exhibiting bias against women.  The article details how Amazon's machine learning system, designed to review job applications, was found to discriminate against female candidates.  The system, trained on resumes submitted over a 10-year period, reflected the male-dominated tech industry, leading to biased recommendations.  The article highlights the potential for bias in AI systems, even when not explicitly programmed with discriminatory criteria.  The context of the course is likely to discuss the ethical implications of AI and the importance of fairness and unbiased algorithms in real-world applications.%
\subsection{Figures}%
\begin{figure}[H]%
\centering%
\begin{subfigure}[b]{0.15\textwidth}%
\centering%
\includegraphics[width=\textwidth]{/Users/ashoksaravanan/Coding/ScribeLec/server/output/CS 243/lectures/16L2-ethics/figures/23.1.png}%
\caption{}%
\end{subfigure}%
\hspace{0.05\textwidth}%
\begin{subfigure}[b]{0.15\textwidth}%
\centering%
\includegraphics[width=\textwidth]{/Users/ashoksaravanan/Coding/ScribeLec/server/output/CS 243/lectures/16L2-ethics/figures/23.2.png}%
\caption{}%
\end{subfigure}%
\hspace{0.05\textwidth}%
\begin{subfigure}[b]{0.15\textwidth}%
\centering%

\includegraphics[width=\textwidth]{/Users/ashoksaravanan/Coding/ScribeLec/server/output/CS 243/lectures/16L2-ethics/figures/23.3.png}%
\caption{}%
\end{subfigure}%
\hspace{0.05\textwidth}%
\begin{subfigure}[b]{0.15\textwidth}%
\centering%
\includegraphics[width=\textwidth]{/Users/ashoksaravanan/Coding/ScribeLec/server/output/CS 243/lectures/16L2-ethics/figures/23.4.png}%
\caption{}%
\end{subfigure}%
\hspace{0.05\textwidth}%
\begin{subfigure}[b]{0.15\textwidth}%
\centering%
\includegraphics[width=\textwidth]{/Users/ashoksaravanan/Coding/ScribeLec/server/output/CS 243/lectures/16L2-ethics/figures/23.5.png}%
\caption{}%
\end{subfigure}%
\end{figure}

%
\newpage%
\section{Slide 24}%
\label{sec:Slide24}%
\subsection{Content}%

\textbf{Potential Solutions: constrain the algorithm}
%
\subsection{Description}%
This slide, likely part of a presentation or lecture on Artificial Intelligence Basics, lists a potential solution to a problem related to algorithms.  The solution is to "constrain the algorithm."  This suggests that the problem involves undesirable behavior or outputs from the algorithm, and the proposed solution is to limit or control the algorithm's actions or choices to achieve a desired outcome.  The context of the course is crucial to understand the specific problem being addressed.  It could be related to bias mitigation, overfitting, or other issues in machine learning algorithms.%
\subsection{Figures}%
\begin{figure}[H]%
\centering%
\end{figure}

%
\newpage%
\section{Slide 25}%
\label{sec:Slide25}%
\subsection{Content}%

\textbf{Fairness through Blindness}

Ignore all irrelevant/protected attributes

But you don't need to see an attribute to be able to predict it with high accuracy

\begin{itemize}
    \item E.g. User visits artofmaniless.com ... has 90\% chance of being male
\end{itemize}
%
\subsection{Description}%
This slide, likely part of a lecture on Artificial Intelligence Basics, discusses a concept called "Fairness through Blindness."  It proposes a method for achieving fairness in machine learning models by ignoring irrelevant or protected attributes.  The key idea is that a model can still achieve high accuracy in prediction without needing to consider potentially discriminatory attributes.  The example given is that a model can predict a high probability of a user being male based on their visit to a website (artofmaniless.com) without needing to know the user's gender.  This suggests a potential trade-off between fairness and accuracy, where ignoring certain attributes might lead to a more fair model but potentially lower accuracy in some cases. The context of the course is likely to discuss the ethical implications of AI and the challenges of building fair and unbiased machine learning models.%
\subsection{Figures}%
\begin{figure}[H]%
\centering%
\end{figure}

%
\newpage%
\section{Slide 26}%
\label{sec:Slide26}%
\subsection{Content}%

\textbf{Group Fairness}

Equalize two groups $S$, $T$ at the level of outcomes

\begin{itemize}
    \item For example, $S$=minority, $T$=everyone else
    \item $P(\text{outcome } o | S) = P(\text{outcome } o | T)$
\end{itemize}

This is not strong enough as a notion of fairness

\begin{itemize}
    \item Sometimes desirable, but can be abused
    \item Self-fulfilling prophecy: select smartest students in $T$, random students in $S$, then students in $T$ will perform better
\end{itemize}
%
\subsection{Description}%
This slide, likely part of a lecture on Artificial Intelligence Basics, discusses the concept of "Group Fairness."  It describes a method of equalizing two groups, $S$ and $T$, at the level of outcomes.  The example provided is comparing a minority group ($S$) to a majority group ($T$).  The key idea is to ensure that the probability of a particular outcome given membership in group $S$ is equal to the probability of that same outcome given membership in group $T$.  Mathematically, this is expressed as $P(\text{outcome } o | S) = P(\text{outcome } o | T)$.  The slide then points out that this approach alone is insufficient for true fairness.  It highlights potential issues like the abuse of such a system and the possibility of self-fulfilling prophecies.  The example of selecting the smartest students for group $T$ and random students for group $S$ illustrates how this seemingly fair approach can lead to worse outcomes for group $S$ due to the inherent biases in the selection process.  The context of the course is likely to be discussing the ethical considerations and limitations of various fairness metrics in AI systems.%
\subsection{Figures}%
\begin{figure}[H]%
\centering%
\end{figure}

%
\newpage%
\section{Slide 27}%
\label{sec:Slide27}%
\subsection{Content}%

\textbf{Individual fairness}

Treat similar individuals similarly

\begin{itemize}
    \item Similar for the purpose of the classification task
    \item Similarly distributed over outcomes
\end{itemize}

How to do this is a current area of research.
%
\subsection{Description}%
This slide, likely part of a lecture on Artificial Intelligence Basics, discusses "Individual Fairness."  It proposes a principle for fair AI systems: treat individuals who are similar in relevant ways similarly.  The specific criteria for similarity are not defined, but the slide implies that similarity should be considered in the context of the classification task itself.  Furthermore, the individuals should be similarly distributed across the possible outcomes.  The slide concludes by stating that developing methods to achieve this principle is an active area of research.  The context of the course suggests a focus on the ethical considerations and practical challenges of building fair and unbiased AI systems.%
\subsection{Figures}%
\begin{figure}[H]%
\centering%
\begin{subfigure}[b]{0.15\textwidth}%
\centering%
\includegraphics[width=\textwidth]{/Users/ashoksaravanan/Coding/ScribeLec/server/output/CS 243/lectures/16L2-ethics/figures/27.1.png}%
\caption{}%
\end{subfigure}%
\end{figure}

%
\newpage%
\section{Slide 28}%
\label{sec:Slide28}%
\subsection{Content}%

\textbf{Interpretable results}
%
\subsection{Description}%
This slide, likely part of a lecture on Artificial Intelligence Basics, presents a single heading: "Interpretable results."  It likely introduces a discussion point about the importance of interpretability in AI models.  The context of the course suggests a focus on the ethical and practical considerations of building AI systems, including the need for models that are not only accurate but also understandable to humans.  The lack of any further content on the slide indicates that this is a title slide or a brief introductory statement for a subsequent discussion.%
\subsection{Figures}%
\begin{figure}[H]%
\centering%
\end{figure}

%
\newpage%
\section{Slide 29}%
\label{sec:Slide29}%
\subsection{Content}%

\textbf{Can we Trust the Output of a ML Algorithm?}

Sometimes you don't know if you can trust a machine learning prediction...

\begin{itemize}
    \item[ ] \textbf{Jane Doe}
    \item[ ] has the flu
    \item[ ] \textbf{Why should I trust you?}
\end{itemize}

Source: Richard Johansson, University of Gothenburg
%
\subsection{Description}%
This slide, likely part of a lecture on Artificial Intelligence Basics, poses a question about the trustworthiness of machine learning (ML) algorithms.  It uses a simple visual representation to illustrate the concept.  A cartoon image of a laptop, displaying a network diagram (likely representing a machine learning model), is shown.  A cartoon character, resembling a doctor, is shown asking "Why should I trust you?".  The text "Jane Doe has the flu" is in a speech bubble next to the laptop, suggesting that the ML algorithm might have predicted that Jane Doe has the flu.  The doctor's question implies a concern about the reliability of the prediction.  The slide highlights the importance of evaluating the output of ML algorithms and the need for transparency and justification in their predictions.  The source is credited as Richard Johansson, University of Gothenburg.%
\subsection{Figures}%
\begin{figure}[H]%
\centering%
\begin{subfigure}[b]{0.15\textwidth}%
\centering%
\includegraphics[width=\textwidth]{/Users/ashoksaravanan/Coding/ScribeLec/server/output/CS 243/lectures/16L2-ethics/figures/29.1.png}%
\caption{}%
\end{subfigure}%
\end{figure}

%
\newpage%
\section{Slide 30}%
\label{sec:Slide30}%
\subsection{Content}%

\textbf{Can We Understand What the Model is Doing?}

It's easier to trust a prediction if you understand the reasons for it.

\begin{itemize}
    \item Fever
    \item Headache
    \item Fatigue
\end{itemize}

Jane Doe
has the flu

Ok. I can trust you.

Source: Richard Johansson, University of Gothenburg
%
\subsection{Description}%
This slide, likely part of a lecture on Artificial Intelligence Basics, discusses the importance of interpretability in machine learning models.  It presents a simple visual example.  A cartoon doctor is shown, and a speech bubble indicates that the doctor is saying "Ok. I can trust you."  The doctor is looking at a laptop with a magnifying glass over a network diagram.  The laptop screen displays the prediction "Jane Doe has the flu."  The slide also shows three icons representing possible symptoms: a thermometer (Fever), a pill bottle (Headache), and a chair (Fatigue).  These icons are in a bubble above the laptop, suggesting that these are the factors considered by the model to arrive at the prediction.  The overall message is that understanding the reasoning behind a machine learning model's prediction increases trust in the output. The source is credited to Richard Johansson, University of Gothenburg.%
\subsection{Figures}%
\begin{figure}[H]%
\centering%
\begin{subfigure}[b]{0.15\textwidth}%
\centering%

\includegraphics[width=\textwidth]{/Users/ashoksaravanan/Coding/ScribeLec/server/output/CS 243/lectures/16L2-ethics/figures/30.1.png}%
\caption{}%
\end{subfigure}%
\end{figure}

%
\newpage%
\section{Slide 31}%
\label{sec:Slide31}%
\subsection{Content}%

\textbf{Understand When the Algorithm Makes No Sense}

Or figuring out when you shouldn't trust a model

\begin{itemize}
    \item You are detecting snow, not wolves!
    \item I can't trust you!
\end{itemize}

\begin{itemize}
    \item I am 100\% sure this is a wolf
\end{itemize}

Source: Richard Johansson, University of Gothenburg
%
\subsection{Description}%
This slide, part of a lecture on Artificial Intelligence Basics, focuses on the importance of evaluating the output of machine learning algorithms.  It presents a visual example of a situation where an algorithm's prediction might be incorrect.  A cartoon image depicts a person looking at a laptop screen.  A speech bubble coming from the laptop says "I am 100% sure this is a wolf."  Another speech bubble coming from the person says "You are detecting snow, not wolves! I can't trust you!"  The image on the laptop screen shows a picture that appears to be a snowy landscape.  The slide highlights the need to critically assess the algorithm's output and identify cases where the prediction is not aligned with reality.  The context of the course suggests a discussion on model evaluation, interpretability, and potential errors in machine learning predictions. The source is credited to Richard Johansson, University of Gothenburg.%
\subsection{Figures}%
\begin{figure}[H]%
\centering%
\begin{subfigure}[b]{0.15\textwidth}%
\centering%
\includegraphics[width=\textwidth]{/Users/ashoksaravanan/Coding/ScribeLec/server/output/CS 243/lectures/16L2-ethics/figures/31.1.png}%
\caption{}%
\end{subfigure}%
\end{figure}

%
\newpage%
\section{Slide 32}%
\label{sec:Slide32}%
\subsection{Content}%

\textbf{Intelligible Models for HealthCare}

Caruana et al. 2015

In the mid 90's, Cost-Effective HealthCare (CEHC)
funded large study to evaluate the application of
machine learning to predicting pneumonia risk

\textbf{Goal:} predict the probability of death (POD) for
patients with pneumonia so high-risk patients could
be admitted to the hospital while low-risk patients
treated as outpatients

Most accurate models were neural nets, which
outperformed traditional methods by wide margin

But after careful consideration they were considered
too risky for use on real patients and logistic
regression was used instead. Why?
%
\subsection{Description}%
This slide, likely from a presentation or lecture on Artificial Intelligence in Healthcare, discusses a study conducted by Caruana et al. in 2015.  The study, funded by Cost-Effective HealthCare (CEHC), investigated the use of machine learning to predict pneumonia risk.  The primary goal was to predict the probability of death (POD) for patients with pneumonia.  The study aimed to identify high-risk patients who should be hospitalized, while low-risk patients could be treated as outpatients.  The slide highlights that neural networks were the most accurate models in the initial stages of the study, significantly outperforming traditional methods.  However, further consideration led to the decision to use logistic regression instead of neural networks due to concerns about the risk associated with using neural networks on real patients.  The slide concludes with a question: "Why?" This implies a discussion point about the reasons behind this decision, potentially exploring factors like model interpretability, explainability, or the potential for errors in neural network predictions in a healthcare context.%
\subsection{Figures}%
\begin{figure}[H]%
\centering%
\begin{subfigure}[b]{0.15\textwidth}%
\centering%
\includegraphics[width=\textwidth]{/Users/ashoksaravanan/Coding/ScribeLec/server/output/CS 243/lectures/16L2-ethics/figures/32.1.png}%
\caption{}%
\end{subfigure}%
\end{figure}

%
\newpage%
\section{Slide 33}%
\label{sec:Slide33}%
\subsection{Content}%

\textbf{Intelligible Models for HealthCare}

Caruana et al. 2015

Rule-based system used in the study learned the
rule \texttt{HasAsthma(x)} $\implies$ \texttt{LowerRisk(x)}

\begin{itemize}
    \item i.e., patients who have a history of asthma have
    lower risk of dying from pneumonia than the
    general population -- which is counterintuitive
\end{itemize}

This reflected a true pattern in the training data

\begin{itemize}
    \item Patients with history of asthma were admitted
    directly to the ICU with aggressive care that
    effectively lowered their risk of dying
\end{itemize}

But causal reasoning is incorrect... asthmatics have
much higher risk if not hospitalized

If rule-based systems have this problem,
then neural networks probably made similar
mistakes... but even more difficult to interpret
%
\subsection{Description}%
This slide discusses a study by Caruana et al. (2015) on using machine learning for healthcare, specifically focusing on the issue of interpretability.  A rule-based system learned a rule that having asthma (\texttt{HasAsthma(x)}) implied a lower risk of death from pneumonia (\texttt{LowerRisk(x)}).  This rule, while seemingly counterintuitive, was a true pattern in the training data.  The slide explains that patients with asthma were admitted to the ICU and received aggressive treatment, which lowered their risk of death.  However, the causal reasoning behind this rule is flawed.  In reality, asthmatics have a much higher risk of death if not hospitalized.  The slide concludes by noting that if rule-based systems exhibit this type of counterintuitive behavior, then neural networks, while potentially even more complex, are likely to exhibit similar issues, making them difficult to interpret. The context of the course is Artificial Intelligence Basics, focusing on the limitations and complexities of machine learning models in real-world applications, particularly in healthcare.%
\subsection{Figures}%
\begin{figure}[H]%
\centering%
\begin{subfigure}[b]{0.15\textwidth}%
\centering%
\includegraphics[width=\textwidth]{/Users/ashoksaravanan/Coding/ScribeLec/server/output/CS 243/lectures/16L2-ethics/figures/33.1.png}%
\caption{}%
\end{subfigure}%
\end{figure}

%
\newpage%
\section{Slide 34}%
\label{sec:Slide34}%
\subsection{Content}%

\textbf{Detecting and removing bias from models}

Bias is in the data, not the leaning algorithm

Models trained on data will lean any biases in the data

\begin{itemize}
    \item ML for resume processing will learn gender bias if data has a gender bias
    \item ML for recidivism prediction will learn race bias if data has a race bias
    \item ...
\end{itemize}

How to deal with bias using transparent models

\begin{itemize}
    \item Must keep bias features in data when model is trained
    \item Remove what was learned from these bias features after training
\end{itemize}

If offending bias variables are eliminated prior to training:

\begin{itemize}
    \item Often can't tell you still have a problem
    \item Makes it harder to correct the problem
\end{itemize}

EU General Data Protection Regulation (in effect 2018)

\begin{itemize}
    \item Article 9 makes it more difficult to use personal data revealing racial or ethnic origin and other "special categories"
\end{itemize}
%
\subsection{Description}%
This slide, from a lecture on Artificial Intelligence Basics, discusses the crucial issue of bias in machine learning models.  It emphasizes that bias originates from the data, not the algorithm itself.  The slide details how models trained on biased data will inevitably reflect those biases.  Examples are given, such as how an ML model for resume processing might learn gender bias if the training data reflects gender bias, and how an ML model for recidivism prediction might learn racial bias if the training data reflects racial bias.  The slide then outlines a strategy for dealing with bias using "transparent models," which involve retaining the biased features in the training data while removing the learned bias after training.  Finally, the slide highlights the challenges of eliminating bias variables prior to training, noting that this approach often fails to identify the presence of bias and can make bias correction more difficult.  The slide also mentions the EU General Data Protection Regulation (GDPR) and how Article 9 of the GDPR makes it harder to use personal data that reveals racial or ethnic origin or other sensitive categories.  The overall message is that addressing bias in machine learning requires careful consideration of the data and the models used, and that simply removing biased features from the data before training is not a sufficient solution.%
\subsection{Figures}%
\begin{figure}[H]%
\centering%
\begin{subfigure}[b]{0.15\textwidth}%
\centering%
\includegraphics[width=\textwidth]{/Users/ashoksaravanan/Coding/ScribeLec/server/output/CS 243/lectures/16L2-ethics/figures/34.1.png}%
\caption{}%
\end{subfigure}%
\end{figure}

%
\newpage%
\section{Slide 35}%
\label{sec:Slide35}%
\subsection{Content}%
%
\subsection{Description}%
This slide is a news article headline and excerpt from Forbes.com, discussing a new New York City law requiring audits for AI biases.  The article, by Lance Eliot, a renowned AI and Machine Learning expert, highlights the potential for problems with the new law, which is set to take effect in January 2023. The article suggests that despite good intentions, the law's design may lead to confusion, consternation, and troubles. The slide includes a title, author, date, and a link to the full article.  The context of the course (Artificial Intelligence Basics) is relevant because it discusses the practical implications and potential challenges of implementing AI regulations in the real world, particularly regarding bias detection and mitigation.%
\subsection{Figures}%
\begin{figure}[H]%
\centering%
\begin{subfigure}[b]{0.15\textwidth}%
\centering%

\includegraphics[width=\textwidth]{/Users/ashoksaravanan/Coding/ScribeLec/server/output/CS 243/lectures/16L2-ethics/figures/35.1.png}%
\caption{}%
\end{subfigure}%
\hspace{0.05\textwidth}%
\begin{subfigure}[b]{0.15\textwidth}%
\centering%
\includegraphics[width=\textwidth]{/Users/ashoksaravanan/Coding/ScribeLec/server/output/CS 243/lectures/16L2-ethics/figures/35.2.png}%
\caption{}%
\end{subfigure}%
\end{figure}

%
\newpage%
\section{Slide 36}%
\label{sec:Slide36}%
\subsection{Content}%

\textbf{Follow-up}

\textbf{NYC's Anti-Bias AI Hiring Law Has Largely Failed So Far}

\textit{Jackie Turner - May 15, 2024}

\textit{Although other jurisdictions have shown interest in passing laws in a similar vein to LL144, lawmakers have been hesitant to move forward with them after seeing NYC's efforts in this area have largely flopped.}

\textit{Last year, New York City became the first U.S. jurisdiction to regulate automated and artificial intelligence (AI) tools in the hiring process. So far, however, employers have paid this law little, if any, attention. The law (Local Law 144) is designed to prevent potential race and gender bias in AI hiring algorithms by requiring employers to regularly audit these tools. However, as employers have the freedom to decide whether or not their systems are covered by the law and overall lax enforcement, it's largely gone ignored. Yet, despite the failure of LL144, employers should still anticipate new and more comprehensive regulation surrounding AI and employment into play in the near future.}

\textit{New York City Passed an AI Hiring Law. So Far, Few Companies Are Following It.}

\textit{By Lauren Weber - Jan. 22, 2024}

\textit{The law requires employers to disclose how algorithms screen job candidates.}

\textit{It has been six months since New York City began enforcing the nation's first law requiring companies to disclose how algorithms influence their hiring decisions. So far, disclosures are rare.}

\textit{The law requires employers that use software to assist with hiring and promotion decisions-including chatbot interviewing tools and résumé scanners-to disclose how algorithms influence their hiring decisions.}

\textit{https://www.wsj.com/business/new-york-city-passed-an-ai-hiring-law-so-far-few-companies-are-following-it-7e31a5b7}
%
\subsection{Description}%
This slide presents news articles discussing the effectiveness of New York City's Local Law 144, a law aimed at regulating AI bias in hiring processes.  The first article, by Jackie Turner, highlights the lack of compliance with the law by employers, suggesting that the law has largely failed to achieve its intended goals.  The second article, by Lauren Weber, further emphasizes this point, noting that six months after the law's implementation, few companies are disclosing how their algorithms influence hiring decisions.  The articles suggest that despite the law's intent to prevent bias in AI-driven hiring, there's been limited adoption and enforcement, leading to a lack of transparency in the hiring process.  The context of the course (Artificial Intelligence Basics) is relevant because it discusses the practical implications and potential challenges of implementing AI regulations in the real world, particularly regarding bias detection and mitigation.  The slide includes links to the full articles.%
\subsection{Figures}%
\begin{figure}[H]%
\centering%
\begin{subfigure}[b]{0.15\textwidth}%
\centering%
\includegraphics[width=\textwidth]{/Users/ashoksaravanan/Coding/ScribeLec/server/output/CS 243/lectures/16L2-ethics/figures/36.1.png}%
\caption{}%
\end{subfigure}%
\hspace{0.05\textwidth}%
\begin{subfigure}[b]{0.15\textwidth}%
\centering%
\includegraphics[width=\textwidth]{/Users/ashoksaravanan/Coding/ScribeLec/server/output/CS 243/lectures/16L2-ethics/figures/36.2.png}%
\caption{}%
\end{subfigure}%
\end{figure}

%
\newpage%
\section{Slide 37}%
\label{sec:Slide37}%
\subsection{Content}%

\textbf{WEAPONS OF} \\
\textbf{MATH DESTRUCTION} \\
HOW BIG DATA INCREASES INEQUALITY \\
AND THREATENS DEMOCRACY \\
CATHY O'NEIL
%
\subsection{Description}%
This slide displays the title and subtitle of a book, "Weapons of Math Destruction," by Cathy O'Neil.  The title is presented in a bold, sans-serif font, and the subtitle is in a smaller, regular font.  The title and subtitle are positioned prominently on a yellow background with red radiating lines, suggesting a dynamic and potentially problematic aspect of the topic.  A stylized image of a skull with interconnected gears and lines is also present, further emphasizing the theme of the book.  The context of the course (Artificial Intelligence Basics) suggests that this slide is likely a reference to the potential negative consequences of using mathematical models and algorithms, particularly in areas like data analysis and machine learning, and how these tools can exacerbate existing societal inequalities.%
\subsection{Figures}%
\begin{figure}[H]%
\centering%
\begin{subfigure}[b]{0.15\textwidth}%
\centering%
\includegraphics[width=\textwidth]{/Users/ashoksaravanan/Coding/ScribeLec/server/output/CS 243/lectures/16L2-ethics/figures/37.1.png}%
\caption{}%
\end{subfigure}%
\end{figure}

%
\newpage%
\end{document}